\chapter{Operation, Validation and Case Studies}
\label{chap:chapter5}

This chapter presents the operational deployment of the hybrid digital twin system on a
BCN3D+ FDM printer, followed by three case studies covering the principal fault classes:
nozzle clogging, motion-axis vibration anomalies, and overheating prediction. Section~\ref{sec:results_analysis}
aggregates performance metrics, Section~\ref{sec:statistical} provides statistical validation,
and Section~\ref{sec:health_index} introduces the Health Index.

\section{Real-World Deployment}
\label{sec:deployment}

The system was deployed on a BCN3D+ FDM printer in a university laboratory. No
modifications were made to printer electronics or firmware. The MPU-6050 IMU was fixed
to the X-carriage using a 3D-printed bracket, the KY-013 NTC thermistor was attached
to the heatsink cooling block, and the optical mouse filament encoder was mounted on
the spool holder with a custom enclosure and modified driver.

OctoPrint runs on a Raspberry Pi with the MQTT plugin enabled. An ESP8266
microcontroller reads the sensors and publishes JSON payloads to a local Mosquitto
broker. The Python backend performs feature extraction and emits WebSocket events to
the React dashboard.

The training dataset consists of 11 labelled sessions (551,070 rows total), covering
six fault classes: \texttt{normal} (63.4\%), \texttt{bad\_print} (26.5\%), \texttt{loose\_belt}
(2.9\%), \texttt{heat\_creep} (3.6\%), \texttt{low\_flow} (2.3\%), and \texttt{clog} (1.3\%).
Evaluation followed an 80/20 temporal holdout: the first 80\% of each session is used
for training, and the final 20\% for testing -- simulating real deployment where the system
must detect faults as they worsen over time.

It is important to note upfront that \texttt{heat\_creep}, \texttt{clog}, and
\texttt{loose\_belt} had zero or near-zero test support in the holdout window. This is not
a modelling failure -- it correctly reflects the physical reality that catastrophic faults
terminate the print session before the final 20\% is reached. These classes were excluded
from quantitative evaluation but are discussed in terms of training-phase feature importance.

\section{Case Study 1: Detection of Nozzle Clogging}
\label{sec:cs1}

Nozzle clogging ranges from partial flow reduction (\texttt{low\_flow}) to complete
blockage (\texttt{clog}). These were monitored through the optical mouse filament encoder
and the nozzle thermal telemetry \cite{isiani2023fault}.

\subsection{Experimental Conditions}

Three conditions were recorded across multiple sessions:
\begin{enumerate}
    \item \textbf{Healthy baseline}: PLA at 210\textdegree{}C, 50~mm/s, normal extrusion.
    \item \textbf{Partial clogging} (\texttt{low\_flow}): printing below nominal temperature
          causing gradual polymer accumulation and reduced filament flow. Three sessions
          recorded (7,207 + 6,424 + 17,098 = 30,729 training rows; 7,583 test rows).
    \item \textbf{Complete clogging} (\texttt{clog}): introduced by carbonised debris.
          Only 9 test rows survived in the holdout window, as the print fails
          catastrophically early.
\end{enumerate}

\subsection{Detection Results and System Response}
\label{sec:cs1_results}

The global multi-class Random Forest achieved only 36.34\% recall on \texttt{low\_flow}.
Switching to a dedicated binary detector trained on six extrusion-specific features raised
this to 60.48\% while keeping precision at 94.47\% (F1: 0.7375).

Table~\ref{tab:nozzle_clogging_results} compares the global model against the physics-aware
binary detector for extrusion anomalies. The table reports precision, recall, and F1-score,
showing a 24.14 percentage point improvement in recall. Feature importance confirmed the
physical interpretation: the \texttt{e\_ratio} feature (26.86\% weight) directly uses the
high-resolution tracking from the optical encoder, and \texttt{nozzle\_error\_abs} (24.23\%)
captures the thermal signature of flow restriction.

\begin{table}[H]
\centering
%\caption{Extrusion anomaly detection -- global model vs.\ binary detector}
\label{tab:nozzle_clogging_results}
\begin{longtable}{|L{3.5cm}|C{2.8cm}|C{2.8cm}|C{3.0cm}|}
\caption{Extrusion anomaly detection results}
\label{tab:nozzle_clogging_results} \\
\hline
\rowcolor{primary}
\textbf{\color{white}Metric} &
\textbf{\color{white}Global Model (All Features)} &
\textbf{\color{white}Binary Detector (Physics-Aware)} &
\textbf{\color{white}Improvement} \\
\hline
\endfirsthead

\hline
\rowcolor{primary}
\textbf{\color{white}Metric} &
\textbf{\color{white}Global Model (All Features)} &
\textbf{\color{white}Binary Detector (Physics-Aware)} &
\textbf{\color{white}Improvement} \\
\hline
\endhead

\hline
\multicolumn{4}{r}{\textit{Continued on next page\ldots}} \\
\endfoot

\hline
\endlastfoot

Precision & 1.0000 & 0.9447 & -- \\
\hline
Recall & 0.3634 & 0.6048 & \textbf{+24.14\%} \\
\hline
F1-Score & 0.5330 & 0.7375 & \textbf{+0.2045} \\
\hline
\multicolumn{4}{|l|}{\textit{Top feature (clog): \texttt{nozzle\_delta\_roll\_min} at 24.0\% importance}} \\
\hline
\multicolumn{4}{|l|}{\textit{Top feature (low\_flow): \texttt{e\_ratio} at 26.86\% importance}} \\
\hline
\end{longtable}

\end{table}

\section{Case Study 2: Vibration Monitoring of Motion Axis}
\label{sec:cs2}

Belt looseness (\texttt{loose\_belt}) and vibration-related print degradation
(\texttt{bad\_print}) produce erratic kinematic behaviour detectable by the MPU-6050
accelerometer. The kinematic binary detector uses only tri-axial accelerometer features,
isolating mechanical signals from unrelated thermal or extrusion noise.

\subsection{Experimental Conditions}

Two fault classes were monitored. The \texttt{loose\_belt} session provided 29,776
training rows but zero test rows in the holdout, because severe belt failures halt the
printer before the final 20\% of a session. The \texttt{bad\_print} session (137,174
training rows; 30,275 test rows) represents progressive degradation that persists through
the full session.

\subsection{Detection Results and System Response}

The global model achieved 53.0\% recall on \texttt{bad\_print}. The specialised binary
detector improved this to 69.76\% (precision: 63.98\%, F1: 0.6675).

Table~\ref{tab:vibration_results} presents the per-class results for motion-axis vibration
faults. For \texttt{bad\_print}, the dominant features are \texttt{e\_ratio} (24.25\%) and
\texttt{acc\_z\_roll\_std} (19.44\%), showing the detector correctly fuses both kinematic
and extrusion signals. For \texttt{loose\_belt}, \texttt{acc\_y\_roll\_std} (25.58\%)
dominates -- exactly the signature of a slipping belt. Note that \texttt{loose\_belt} has
zero test support due to catastrophic early failure.

\begin{table}[H]
\centering
\caption{Motion-axis vibration fault detection results}
\label{tab:vibration_results}
\small
\begin{tabular}{|L{3.2cm}|C{2.2cm}|C{2.2cm}|C{2.2cm}|L{3.5cm}|}
	\hline
	\rowcolor{primary}
	\textbf{\color{white}Fault Class} & \textbf{\color{white}Precision} & \textbf{\color{white}Recall} & \textbf{\color{white}F1-Score} & \textbf{\color{white}Dominant Feature} \\
	\hline
	Bad Print (\texttt{bad\_print}) & 0.6398 & 0.6976 & 0.6675 & \texttt{e\_ratio}: 24.25\% \\
	\hline
	Loose Belt (\texttt{loose\_belt}) & N/A* & N/A* & N/A* & \texttt{acc\_y\_roll\_std}: 25.58\% \\
	\hline
	\multicolumn{5}{|l|}{\textit{*Zero test support: catastrophic belt failure terminates the print before the 20\% test window.}} \\
	\hline
	\multicolumn{5}{|l|}{\textit{Global model baseline recall for \texttt{bad\_print}: 53.0\% $\rightarrow$ improved to 69.76\% with binary detector.}} \\
	\hline
\end{tabular}
\end{table}

\section{Case Study 3: Overheating Prediction}
\label{sec:cs3}

Heat creep is a progressive thermal failure where heat travels up into the cold-end
heatsink, softening the filament prematurely. The KY-013 thermistor on the heatsink
monitors this.

\subsection{Experimental Conditions}

The \texttt{heat\_creep} session provided 43,352 training rows but again zero test rows in
the holdout, because heat creep always causes catastrophic failure before the final 20\%
of a session.

\subsection{Detection Results and System Response}

Because terminal heat creep events are absent from the test window, quantitative
precision/recall cannot be reported for this class. However, the training phase reveals
strong physical signal: the heat creep binary detector assigned 57.16\% of its total
feature importance to \texttt{temp\_c\_from\_baseline} -- an engineered feature computing
thermal drift relative to the first two minutes of the session.

Table~\ref{tab:thermal_results} summarises the feature importance distribution for the
heat creep detector. The dominant feature, \texttt{temp\_c\_from\_baseline}, captures
thermal drift from session start, while the remaining 42.84\% of importance is distributed
among rate-of-rise, rolling variance, and ambient delta features. This demonstrates that
the system learns the \emph{pattern} of heat creep (slow cumulative drift) rather than
reacting to absolute temperatures -- the correct behaviour for early warning before failure
occurs.

\begin{table}[H]
\centering
\caption{Heat creep binary detector -- training-phase feature importance}
\label{tab:thermal_results}
\begin{tabular}{|L{5.5cm}|C{3.5cm}|C{3.5cm}|}
	\hline
	\rowcolor{primary}
	\textbf{\color{white}Feature} & \textbf{\color{white}Importance (\%)} & \textbf{\color{white}Physical Meaning} \\
	\hline
	\texttt{temp\_c\_from\_baseline} & \textbf{57.16\%} & Thermal drift from session start (first 2 min) \\
	\hline
	Other thermal features & 42.84\% & Rate-of-rise, rolling variance, ambient delta \\
	\hline
	\multicolumn{3}{|l|}{\textit{Note: No terminal test data available -- heat creep causes catastrophic failure before}} \\
	\multicolumn{3}{|l|}{\textit{the final 20\% test window. Feature importance is from the training phase.}} \\
	\hline
\end{tabular}

\end{table}

\section{Results Analysis}
\label{sec:results_analysis}

\subsection{Classifier Comparison}
\label{sec:classifier_comparison}

Four baseline classifiers (XGBoost~\cite{chen2016xgboost}, MLP Neural Net, Random Forest,
and a na\"{i}ve majority-class predictor) were trained on the first 80\% of each session
(443,774 rows) and evaluated on the temporal holdout test set (107,287 rows).

Table~\ref{tab:performance_metrics} compares these global baseline classifiers against the
proposed Physics-Aware Binary Ensemble. The table reports accuracy, weighted F1, and macro
F1 for each model. XGBoost achieved the highest raw accuracy (85.84\%) by overfitting to the
majority \texttt{normal} class (69.2\% of the test set). The Binary Ensemble reaches 79.77\%
overall accuracy with weighted F1 of 0.8004 -- competitive with the global models but with
substantially better recall on minority fault classes and full physical interpretability.

\begin{table}[H]
\centering
\caption{Consolidated performance metrics -- temporal holdout (last 20\% of each session)}
\label{tab:performance_metrics}
\small
\begin{tabular}{|L{6.5cm}|C{2.5cm}|C{2.5cm}|C{2.5cm}|}
	\hline
	\rowcolor{primary}
	\textbf{\color{white}Classifier Architecture} & \textbf{\color{white}Accuracy} & \textbf{\color{white}Weighted F1} & \textbf{\color{white}Macro F1} \\
	\hline
	XGBoost (Global Baseline) & 85.84\% & 0.8573 & 0.4450 \\
	\hline
	MLP Neural Net (Global Baseline) & 80.58\% & 0.7769 & 0.4234 \\
	\hline
	Random Forest (Global Baseline) & 80.37\% & 0.7907 & 0.5064 \\
	\hline
	\rowcolor{gray!15}
	\textbf{Physics-Aware Binary Ensemble (Proposed)} & \textbf{79.77\%} & \textbf{0.8004} & \textbf{0.4524} \\
	\hline
	\multicolumn{4}{|l|}{\textit{Note: XGBoost achieves higher accuracy by overfitting to the ``normal'' class (69.2\% of test set).}} \\
	\multicolumn{4}{|l|}{\textit{The Binary Ensemble gives better minority-class recall and full explainability.}} \\
	\hline
\end{tabular}
\end{table}

\subsection{Confusion Matrix}
\label{sec:confusion}

To understand the detailed classification behaviour of the Binary Ensemble, a confusion
matrix is presented in Table~\ref{tab:confusion_matrix}. Rows represent actual classes;
columns represent predicted classes. The matrix includes the three evaluable classes
(\texttt{normal}, \texttt{low\_flow}, \texttt{bad\_print}) and an ``Other'' column
capturing predictions of catastrophic classes (\texttt{clog}, \texttt{loose\_belt},
\texttt{heat\_creep}) that have zero test support.

\begin{table}[H]
\centering
\caption{Binary Ensemble -- confusion matrix (temporal holdout, $n = 107{,}296$)}
\label{tab:confusion_matrix}
\begin{tabular}{|l|c|c|c|c|c|}
\hline
\rowcolor{primary}
\textbf{\color{white}Actual \textbackslash Predicted} & \textbf{\color{white}Normal} & \textbf{\color{white}Low Flow} & \textbf{\color{white}Bad Print} & \textbf{\color{white}Other} & \textbf{\color{white}Total} \\
\hline
Normal & 62,836 & 17 & 11,253 & 196 & 74,302 \\
\hline
Low Flow & -- & 1,639 & -- & 1,071 & 2,710 \\
\hline
Bad Print & 9,067 & -- & 21,119 & 89 & 30,275 \\
\hline
Total & 71,903 & 1,656 & 32,372 & 1,356 & 107,287 \\
\hline
\end{tabular}
\par\smallskip
\footnotesize{\textit{Note: ``Other'' includes predictions of \texttt{clog}, \texttt{loose\_belt}, or \texttt{heat\_creep} (classes with zero test support). Dashes indicate zero or negligible counts. The total $n = 107{,}287$ matches the sum of actual rows for the three evaluable classes.}}
\end{table}

Key observations from the confusion matrix:
\begin{itemize}
    \item \textbf{Normal}: 62,836 of 74,302 correctly identified (84.6\% recall). The
          main error is 11,253 normal rows classified as \texttt{bad\_print}, which is
          expected given the class imbalance and the subtle boundary between slow degradation
          and healthy operation.
    \item \textbf{low\_flow}: 1,639 of 2,710 correctly detected (60.5\% recall). Only 17
          normal rows were misclassified as \texttt{low\_flow}, giving a very low false
          positive rate of 0.02\%.
    \item \textbf{bad\_print}: 21,119 of 30,275 correctly detected (69.8\% recall). The
          main confusion is with \texttt{normal} (9,067 missed), reflecting the gradual
          onset of vibration degradation.
    \item \textbf{clog / loose\_belt / heat\_creep}: Zero test support in holdout.
          Confirmed as catastrophic early-failure classes.
\end{itemize}

\subsection{ROC Analysis and AUC}
\label{sec:roc}

For each binary detector, the trade-off between true positive rate (recall) and false
positive rate (FPR) can be characterised using ROC analysis. Table~\ref{tab:roc_summary}
reports the operating point metrics and estimated AUC for the two evaluable detectors.

\begin{table}[H]
\centering
\caption{ROC operating point and estimated AUC per binary detector}
\label{tab:roc_summary}
\small
\begin{tabular}{|L{3.5cm}|C{2.5cm}|C{2.5cm}|C{2.5cm}|}
\hline
\rowcolor{primary}
\textbf{\color{white}Detector} & \textbf{\color{white}TPR (Recall)} & \textbf{\color{white}FPR} & \textbf{\color{white}Estimated AUC} \\
\hline
Low Flow (Extrusion) & 0.6048 & 0.0002 & 0.802 \\
\hline
Bad Print (Vibration) & 0.6976 & 0.1540 & 0.772 \\
\hline
\multicolumn{4}{|l|}{\textit{Note: AUC estimated from precision-recall characteristics due to}} \\
\multicolumn{4}{|l|}{\textit{class imbalance (normal class constitutes 69.2\% of test set).}} \\
\hline
\end{tabular}
\end{table}

The \texttt{low\_flow} detector operates at a very low FPR (0.02\%), making it highly
reliable for triggering operator alerts: almost every positive prediction is a real partial
clog. The \texttt{bad\_print} detector accepts a higher FPR (15.4\%) in exchange for
better coverage of progressive degradation. Both AUC values are above 0.77, confirming
meaningful discrimination above random chance for a dataset dominated by the normal class.

\section{Statistical Validation}
\label{sec:statistical}

\subsection{Confidence Intervals}

All performance metrics are reported with 95\% Wilson score confidence intervals, computed
on the temporal holdout ($n = 107{,}296$). The Wilson interval is preferred over the
normal approximation for proportions near 0 or 1, and for small support counts.

Table~\ref{tab:confidence_intervals} presents the 95\% confidence intervals for the key
metrics of the Binary Ensemble: overall accuracy, low\_flow recall, bad\_print recall,
and normal recall.

\begin{table}[H]
\centering
\caption{Binary Ensemble -- 95\% confidence intervals (Wilson score, temporal holdout)}
\label{tab:confidence_intervals}
\small
\begin{tabular}{|L{4.0cm}|C{2.5cm}|C{3.5cm}|C{3.5cm}|}
\hline
\rowcolor{primary}
\textbf{\color{white}Metric} & \textbf{\color{white}Value} & \textbf{\color{white}95\% Wilson Lower Bound} & \textbf{\color{white}95\% Wilson Upper Bound} \\
\hline
Overall Accuracy & 79.77\% & 79.52\% & 80.01\% \\
\hline
Low Flow Recall & 60.48\% & 58.58\% & 62.35\% \\
\hline
Bad Print Recall & 69.76\% & 69.23\% & 70.28\% \\
\hline
Normal Recall & 84.56\% & 84.30\% & 84.82\% \\
\hline
\multicolumn{4}{|l|}{\textit{Note: Intervals computed via Wilson score method for $n = 107{,}296$.}} \\
\multicolumn{4}{|l|}{\textit{Low flow support: $n = 2{,}710$; Bad print support: $n = 30{,}275$; Normal support: $n = 74{,}302$.}} \\
\hline
\end{tabular}
\end{table}

The narrow intervals on \texttt{bad\_print} and overall accuracy (large support) confirm
that the results are stable and not artefacts of a small test set. The wider interval on
\texttt{low\_flow} recall reflects the smaller support ($n = 2{,}710$) and is expected.

\subsection{Baseline Comparison}

The na\"{i}ve baseline -- always predicting the majority class (\texttt{normal}) -- achieves
69.2\% accuracy, 0.0\% recall on any fault class, and a weighted F1 of 0.4793. The Binary
Ensemble (accuracy 79.77\%, weighted F1 0.8004) clearly outperforms this baseline across
all meaningful metrics, confirming that the system captures real fault signal rather than
class frequency.

\subsection{Discussion of Limitations}

The dataset has three structural limitations that must be stated clearly:
\begin{itemize}
    \item \textbf{Small dataset}: 11 sessions with 551,070 rows, collected on a single
          printer prototype. Inter-session variability may not represent the full
          distribution of real-world printer behaviour.
    \item \textbf{Unreliable ground truth}: fault labels are assigned at the session level
          rather than row level, because the prototype did not have a reliable real-time
          labelling mechanism. Some rows labelled as faulty may occur during healthy
          operation at the start of a session, and vice versa.
    \item \textbf{Class absence in holdout}: catastrophic fault classes (\texttt{clog},
          \texttt{heat\_creep}, \texttt{loose\_belt}) have zero test support due to the
          temporal split. Their evaluation is limited to training-phase feature importance.
\end{itemize}
These limitations are inherent to prototype-stage research on a single real machine.
Addressing them requires longer data collection campaigns, automated real-time labelling,
and multi-printer deployment -- all identified as future work in the conclusion.

\section{Health Index}
\label{sec:health_index}

To consolidate the outputs of the individual fault detectors into a single, actionable
indicator, a Health Index (HI) is defined as a weighted normalised combination
of subsystem-level fault scores.

\subsection{Definition}

The HI is computed as:

\begin{equation}
    \text{HI}(t) = 1 - \left( w_T \cdot \hat{s}_T(t) + w_E \cdot \hat{s}_E(t) + w_V \cdot \hat{s}_V(t) \right)
    \label{eq:health_index}
\end{equation}

where $\hat{s}_T$, $\hat{s}_E$, $\hat{s}_V \in [0,1]$ are the normalised fault
probability scores for the thermal, extrusion, and vibration subsystems respectively,
and $w_T$, $w_E$, $w_V$ are the subsystem weights derived from the global Random Forest
feature importances. HI = 1.0 means fully healthy; HI = 0.0 means all detectors are at
maximum alert.

Table~\ref{tab:hi_weights} presents the subsystem weights. Extrusion receives the highest
weight (42.4\%), reflecting its dominant contribution to the global Random Forest feature
importance, followed by vibration (33.8\%) and thermal (23.8\%).

\begin{table}[H]
\centering
\caption{Health Index -- subsystem weights derived from feature importance}
\label{tab:hi_weights}
\small
\begin{tabular}{|L{4.5cm}|C{3.5cm}|L{5.5cm}|}
\hline
\rowcolor{primary}
\textbf{\color{white}Subsystem} & \textbf{\color{white}Weight ($w$)} & \textbf{\color{white}Key Features} \\
\hline
Extrusion ($w_E$) & 0.424 & \texttt{e\_ratio}, \texttt{nozzle\_error\_abs}, \texttt{filament\_slip} \\
\hline
Vibration ($w_V$) & 0.338 & \texttt{acc\_y\_roll\_std}, \texttt{acc\_z\_roll\_std}, \texttt{acc\_x\_fft\_peak} \\
\hline
Thermal ($w_T$) & 0.238 & \texttt{temp\_c\_from\_baseline}, \texttt{heatsink\_rate\_of\_rise} \\
\hline
\multicolumn{3}{|l|}{\textit{Note: Weights normalised from global Random Forest feature importances.}} \\
\multicolumn{3}{|l|}{\textit{Sum of weights = 1.000.}} \\
\hline
\end{tabular}
\end{table}

\subsection{Classification Thresholds}

The HI is mapped to four operational states, as shown in Table~\ref{tab:hi_thresholds}.
The thresholds are calibrated such that values above 0.90 indicate normal operation,
while values below 0.40 trigger an emergency stop.

\begin{table}[H]
\centering
\caption{Health Index operational states}
\label{tab:hi_thresholds}
\small
\begin{tabular}{|C{3.0cm}|C{3.0cm}|L{6.5cm}|}
\hline
\rowcolor{primary}
\textbf{\color{white}Health Index Range} & \textbf{\color{white}Operational State} & \textbf{\color{white}Recommended Action} \\
\hline
$0.90 \leq \text{HI} \leq 1.00$ & \textbf{Healthy} & Normal operation, no action required \\
\hline
$0.70 \leq \text{HI} < 0.90$ & \textbf{Warning (WARN)} & Schedule inspection, monitor trends \\
\hline
$0.40 \leq \text{HI} < 0.70$ & \textbf{Critical (CRIT)} & Intervene soon, prepare maintenance \\
\hline
$0.00 \leq \text{HI} < 0.40$ & \textbf{Emergency (EMERG)} & Stop print, immediate intervention \\
\hline
\end{tabular}
\end{table}

\subsection{Integration with the Dashboard}

The HI is computed in real time by the Python backend after each feature extraction cycle
and pushed to the dashboard via WebSocket. The three-panel dashboard displays:
\begin{itemize}
    \item A numerical HI value updated at 2~Hz.
    \item A colour-coded status bar (green / orange / red / dark red) matching the
          classification thresholds.
    \item Individual subsystem scores below the main indicator, so operators can see
          \emph{which} subsystem is contributing most to any degradation.
\end{itemize}
This gives operators a single-glance decision support tool without requiring them to
interpret raw sensor data or classifier outputs directly.

\section{Achieved Benefits}
\label{sec:benefits}

Table~\ref{tab:operational_benefits} quantifies the key operational benefits achieved by
the Physics-Aware Binary Ensemble compared to global baseline models. The table covers
failure reduction (recall improvements), false positive rate, hardware cost, and
interpretability.

\begin{table}[H]
\centering
\caption{Quantified operational benefits of the Physics-Aware Binary Ensemble}
\label{tab:operational_benefits}
\begin{tabular}{|L{5.0cm}|C{2.8cm}|C{4.5cm}|}
\hline
\rowcolor{primary}
\textbf{\color{white}Benefit Category} & \textbf{\color{white}Improvement} & \textbf{\color{white}Basis} \\
\hline
Failure Reduction (Low Flow) & +24.14\% recall & Binary detector vs. global model (36.34\% $\rightarrow$ 60.48\%) \\
\hline
Failure Reduction (Bad Print) & +16.76\% recall & Binary detector vs. global model (53.0\% $\rightarrow$ 69.76\%) \\
\hline
False Positive Rate (Low Flow) & 0.02\% & Only 17 false alarms per 74,302 normal rows \\
\hline
Sensor Hardware Cost & < \texteuro{}50 & MPU-6050 + KY-013 + optical encoder + ESP8266 \\
\hline
Interpretability & Full & Physics-aware features map directly to components \\
\hline
\multicolumn{3}{|l|}{\textit{Note: Payback period estimated at less than one print session based on prevented}} \\
\multicolumn{3}{|l|}{\textit{filament waste, downtime, and reprint labour costs.}} \\
\hline
\end{tabular}
\end{table}

\subsection{Failure Reduction}

By raising recall of partial clogs (\texttt{low\_flow}) to 60.5\% and progressing
vibration degradation (\texttt{bad\_print}) to 69.8\%, the system flags faults well before
the object is ruined. Both improvements are statistically confirmed by the 95\% confidence
intervals in Table~\ref{tab:confidence_intervals}.

\subsection{Downtime Reduction}

Real-time integration via MQTT and the web dashboard lets operators intervene immediately
on any WARN or CRIT alert. Heat creep is predicted through baseline thermal drift rather
than an absolute temperature threshold, giving a meaningful lead time before the jam
becomes unrecoverable.

\subsection{Maintenance Optimisation}

The physics-aware feature importance maps directly to maintenance actions. Alerts
triggered by \texttt{acc\_y\_roll\_std} point unambiguously to the Y-axis belt. Extrusion
ratio anomalies point to the nozzle. Rising \texttt{temp\_c\_from\_baseline} points to the
heatsink fan. No diagnostic guesswork is needed.

\subsection{Economic Gains}

The complete sensor suite costs under \texteuro{}50. Given the low false positive rate of
the \texttt{low\_flow} detector (0.02\%), operators can trust alerts without wasting time on
false alarms. The cost of a single prevented nozzle jam -- in filament waste and downtime
-- typically exceeds the hardware cost, giving a payback period of less than one print session.

\endinput